\section{Conclusion}

Les trois parties du laboratoire ont été réalisées, simulées et, pour le chenillard, testées avec succès sur la carte MAX\_V\_80P\_25P.

\begin{itemize}
	\item \textbf{Partie 1 :} le registre 4 bits décale à droite et conserve sa valeur lorsque \texttt{en\_sh\_i} = 0, grâce au multiplexeur placé devant chaque bascule.
	\item \textbf{Partie 2 :} la table des fonctions synchrones se traduit directement en un multiplexeur d'état futur. La simulation confirme chaque fonction, la priorité du chargement et le reset asynchrone.
	\item \textbf{Partie 3 :} le chenillard réutilise le registre de la partie 2 sans nouvelle logique : reboucler les bits extrêmes sur les entrées série suffit à transformer le décalage en rotation cyclique. Les résultats sur la carte sont identiques à la simulation, et l'intégration a été validée par M. Messerli.
\end{itemize}

Ce laboratoire montre l'intérêt de la décomposition état futur / registre et de la réutilisation : un composant bien conçu et testé sert de base à un système plus complexe avec très peu de modifications.
